% COMPRESS EDIT (length): §4 cut from ~5,000+ words to a lean core. Four redundant
% main-text floats (tab:three_classifier, tab:opportunity_permutation,
% tab:scope_matrix, tab:validation_benchmarks) and three timing figures
% (fig:opportunity_permutation, fig:rolling_origin_pr_auc, fig:leave_one_case_out)
% demoted to Appendix D (\ref{app:validation_audits_submission}); refs repointed in prose.
\section{Validating the Loser-Side Ranking}
\label{sec:validation}

This section validates the award-layer ranking, not the legal
status of ranked firms. The validation target is the one defined
in Section~\ref{sec:cade_adjacency_labels}: always-loser firms
observed alongside direct CADE defendants in adjudicated
procurement-cartel environments. These cobidders are
adjudication-anchored exposure labels, not legal members of a
cartel. A useful ranking must do more than report a high AUC. The
main threat is exposure---firms that participate often, or operate
in the same environments as defendants, have more opportunities to
appear near those anchors even under ordinary competition, so a
screen must concentrate loser-side firms beyond what participation
volume and opportunity alone imply
\citep{porter1993detection,bajari2003deciding,conley2016detecting,kawai2022detecting}.
Two further threats are timing (the screen should rank before the
target information is used) and legal overreach (it should not be
read as a generic cartel-membership detector).

Subsection~\ref{sec:label_funnel} fixes the label-construction
funnel and reconciles the validation samples.
Subsection~\ref{sec:validation_exposure} is the central exercise:
it disciplines the ranking against procurement opportunity and
isolates the score's residual contribution net of exposure.
Subsection~\ref{sec:validation_timing} imposes strict timing and
case-composition discipline, with the sobering result that the
operational ranking is a retrospective, incumbent-pool, largely
single-case diagnostic.
Subsection~\ref{sec:validation_direct} folds the direct-defendant
scope check into a bounded reading of what the validation
establishes.

\subsection{Label Construction and Sample Reconciliation}
\label{sec:label_funnel}

The full CADE legal portfolio comprises $12$ adjudicated
procurement-cartel cases with $65$ legal firm-defendants. We match
defendants to BEC participant records by exact $14$-digit CNPJ
equality, with no fuzzy or manual step; $\valDirectCADE$ are active
in BEC and define the direct-defendant set. Their defendant
tender-items ($\valFunnelDefItems$ distinct items) locate the
cartel-affected environments, within which a cobidder is a zero-win
firm sharing at least one tender-item with a direct defendant.

The main validation target is the $\valCobidders$ always-loser
frequent-loser cobidders defined under a narrow cartel-tender
restriction, which we keep as the primary benchmark for
comparability with prior drafts. We disclose one replication caveat
cleanly: that target's original builder is archived and is not
reproduced by current scripts. A fully scripted alternative
---counting any always-loser sharing a tender-item with a BEC-active
direct defendant over the full portfolio---yields
$\valFunnelFLcobBroad$ frequent-loser and $\valFunnelALcobBroad$
always-loser cobidders, overlapping the main target in
$\valFunnelStaticOverlap$ firms; we carry it as the reproducible
robustness funnel. The conservative pre-$\valConservativeCutoff$
benchmark is not a mechanical subset of the main target: it
restricts the portfolio to the $\valConservativeCases$ early cases
but applies the broad cobidder definition over the wider
always-loser stratum, yielding $\valConservativeCobidders$
cobidders ($\valConservativeFL$ frequent losers) from
$\valConservativeFD$ defendants. Under a common definition the
subset relation is restored. Table~\ref{tab:label_funnel_submission}
reports every benchmark side by side.

\begin{table}[htbp]
\centering
\caption{Label-Construction Funnel and Benchmark Reconciliation}
\label{tab:label_funnel_submission}
\scriptsize
\begin{adjustbox}{max width=\textwidth,center}
\begin{threeparttable}
\begin{tabular}{p{3.0cm}cccccp{1.7cm}p{2.4cm}}
\toprule
Sample / benchmark & Cases & Legal def. & BEC-active def. & AL cobid. & FL cobid. & Cobidder def. & Reason vs main target \\
\midrule
Full CADE portfolio & 12 & 65 & $\valDirectCADE$ & $\valFunnelALcobBroad$ & $\valFunnelFLcobBroad$ & broad & reference \\
Main target --- static narrow & 12 & 65 & -- & -- & $\valCobidders$ & narrow & primary \\
Transparent broad funnel & 12 & 65 & $\valFunnelDirectActiveFTM$ & $\valFunnelALcobBroad$ & $\valFunnelFLcobBroad$ & broad & robustness \\
Conservative pre-$\valConservativeCutoff$ & 4 & -- & $\valConservativeFD$ & $\valConservativeCobidders$ & $\valConservativeFL$ & broad & early-cases subset \\
\bottomrule
\end{tabular}
\begin{tablenotes}
\scriptsize
\item \textit{Notes:} A cobidder is an adjudication-anchored
exposure label, not membership. The main target ($\valCobidders$,
narrow cartel-tender definition) has an archived builder; the
broad shared-tender-item funnel yields $\valFunnelFLcobBroad$
frequent-loser and $\valFunnelALcobBroad$ always-loser cobidders.
The conservative benchmark reproduces to
$\valConservativeCobidders$/$\valConservativeFL$. The
$\valDirectCADE$ cited active defendants correspond to $48$
distinct crossmatch CNPJ, of which $\valFunnelDirectActiveFTM$ are
firm\_tender\_map-active. Counts are unique firms.
\end{tablenotes}
\end{threeparttable}
\end{adjustbox}
\end{table}

Before adding stricter discipline, a baseline question: do
firms flagged by the award-layer rule appear disproportionately
among always-loser firms near adjudicated anchors? They do. In the
conservative pre-$\valConservativeCutoff$ benchmark, the observed
frequent-loser share among cobidders is $\valCADEenrich$ against a
participation-stratified placebo mean of $\valCADEbase$ (excess
ratio $\valMultThreeTwo\!\times$, $p<0.001$), with ROC-AUC
$\valAUCprePost$ (95\% DeLong CI $\valAUCprePostCI$)
\citep{delong1988comparing}. Persistent zero-win participation is
not random with respect to adjudication-anchored loser-side
exposure. But a high AUC here is not yet enough, because
participation volume and shared environments mechanically
manufacture proximity. The next subsection therefore asks how much
of the ranking survives once opportunity is held comparable.

\subsection{Opportunity-Adjusted Validation}
\label{sec:validation_exposure}
\label{sec:validation_exposure_timing}

This is the central validation exercise. The threat is mechanical:
a firm with high participation $T_i$ co-appears with defendants more
often by construction, and a firm in the same item, buyer, year,
and modality cells is more exposed even under ordinary competition.
We measure opportunity directly. For each always-loser firm we count
observed contact with direct defendants, $O_i$, and construct
expected contact under random matching, $E_i=\sum_g n_{ig}\,p_g$,
where $n_{ig}$ is participation in cell $g$ and $p_g$ its
defendant-density. We build $E_i$ at three cell granularities and
restrict to common support; the exposed subsample retains
$n=\valExpExposedN$ always-loser firms containing all $\valExpNpos$
positives (construction in \ref{app:market_opportunity_submission}).
Holding $E_i$ fixed, we ask whether the score reproduces the label
better than exposure alone, whether it adds anything net of
exposure, and whether it ranks within equal-opportunity strata.

The honest answer has two parts. First, exposure is a strong
baseline. An exposure-only model---no score---reproduces the label
at ROC-AUC $\valExpOnlyAUC$ in the exposed subsample, \emph{above}
the raw score's $\valExpRawScoreAUC$, and its PR-AUC
($\valExpExpOnlyPRAUC$) also exceeds the raw score's
($\valExpRawScorePRAUC$). Combined exposure-plus-score logits reach
ROC-AUC near $0.98$ (Table~\ref{tab:opportunity_adjusted_validation}),
but those numbers are mechanically inflated because $E_i$ is built
from the cells where defendants live and so encodes the label
directly; we report them only to make the inflation visible.

Second, conditional on opportunity the score retains limited but
significant residual value. A nested model that adds the score to
the exposure-only model raises AUC by $\valExpIncrement$ (DeLong
$p=\valExpDeLongP$)---the score's isolated contribution, not the
inflated combined fit. And \emph{within} opportunity strata, where
firms are compared only against others with comparable expected
contact, the continuous score still ranks the label at AUC
$\valExpWithinAUC$ (FL14 $\valExpWithinFLAUC$): well above chance
and far below the raw $\valExpRawScoreAUC$. The score also ranks
positive excess contact $O_i-E_i>0$ at AUC
$\valExpScoreExcessAUC$, and a coarsened-exact opportunity match
yields AUC $\valExpCEMmatchedAUC$ with a within-stratum gap of
$\valExpMatchedDelta$ in cobidder probability between flagged and
non-flagged always-losers.
Figure~\ref{fig:observed_vs_expected_contact} shows the positive
$O_i-E_i$ gap in the top score bins---the within-opportunity signal
made visible.

Two exposure-preserving permutation tests close the case, in
opposite directions. Under a pure-exposure binomial null the
observed PR-AUC $\valExpRawScorePRAUC$ does \emph{not} beat the
null (empirical $p=\valExpPermBp$): a pure-exposure model
out-predicts the raw score on the rare-target metric. Under a
permutation that shuffles cobidder labels only within matched
exposure-by-participation strata, the observed PR-AUC \emph{exceeds}
the null ($p=\valExpPermCp$): once opportunity is matched, the
labels are non-randomly score-ranked, and the frequent-loser
enrichment beats both nulls ($p<\valExpFLenrichP$). Both
permutation distributions are reported in
Appendix~\ref{app:validation_audits_submission}.

The conservative reading is the one we carry forward. The
award-layer score does not reproduce the adjudication-anchored
label any better than firms' procurement opportunity does. What it
adds, net of opportunity, is a small and statistically reliable
residual ordering---triage value, a reason to look first at one
loser-side firm rather than another at the bid-layer stage. It is
not proof of collusion and not a causal claim. The remaining
threats---timing, leakage, and case composition---are taken up next:
a strict-timing battery, a leakage audit holding cobidders out of
score construction (out-of-fold AUC $\valLeakCVAUC$ against an
in-sample reference $\valLeakRefAUC$), and a leave-one-case-out
check motivated by one composition fact: the single largest CADE
case supplies roughly $\valExpTopCaseTPshare$ of true positives at
the top $500$ of the ranking.

\begin{table}[htbp]
\centering
\caption{Opportunity-Adjusted Validation of the Award-Layer Ranking}
\label{tab:opportunity_adjusted_validation}
\scriptsize
\begin{adjustbox}{max width=\textwidth,center}
\begin{threeparttable}
\begin{tabular}{p{3.6cm}lccccccccl}
\toprule
Design & Score & $N$ & Pos. & ROC-AUC & PR-AUC & Prec@500 & Rec@500 & Lift@500 & 95\% CI / $p$ \\
\midrule
Raw score & log\_tc & $16{,}843$ & $\valExpNpos$ & $\valExpRawScoreAUC$ & $\valExpRawScorePRAUC$ & $0.130$ & $0.340$ & $11.46$ & $[0.931,0.946]$ \\
Raw FL14 & fl14 & $16{,}843$ & $\valExpNpos$ & $0.924$ & $\valExpRawFLPRAUC$ & $0.088$ & $0.230$ & $7.76$ & $[0.921,0.926]$ \\
Exposure-only benchmark & $\log E_i$ & $\valExpExposedN$ & $\valExpNpos$ & $\valExpOnlyAUC$ & $\valExpExpOnlyPRAUC$ & $0.182$ & $0.476$ & $4.05$ & $[0.818,0.875]$ \\
Within-opportunity stratum & log\_tc$\,|\,$opp & $\valExpExposedN$ & $\valExpNpos$ & $\valExpWithinAUC$ & --- & --- & --- & --- & pooled $C$ \\
Nested increment over exposure & $+$log\_tc & $\valExpExposedN$ & $\valExpNpos$ & $+\valExpIncrement$ & --- & --- & --- & --- & DeLong $p=\valExpDeLongP$ \\
Score ranks excess contact & log\_tc & $16{,}843$ & $475$ & $\valExpScoreExcessAUC$ & $0.095$ & $0.150$ & $0.158$ & $5.32$ & $1[O_i{-}E_i{>}0]$ \\
Permutation (exposure-preserving) & log\_tc & $16{,}843$ & $\valExpNpos$ & --- & $\valExpRawScorePRAUC$ & --- & --- & --- & $p=\valExpPermBp$ (B) \\
Matched / stratified (CEM-opp) & log\_tc & $\valExpExposedN$ & $\valExpNpos$ & $\valExpCEMmatchedAUC$ & --- & --- & --- & --- & $\Delta=\valExpMatchedDelta$ \\
\bottomrule
\end{tabular}
\begin{tablenotes}
\scriptsize
\item \textit{Notes:} The validation target is
adjudication-anchored exposure (cobidder status) among
always-loser firms, not cartel membership. Expected contact
$E_i=\sum_g n_{ig}p_g$ is built from procurement opportunity
cells (item, buyer, year, modality); cell definitions and
retained support are in \ref{app:market_opportunity_submission}.
PR-AUC is the primary metric because the target is rare
($\valExpNpos/16{,}843$); ROC-AUC is reported alongside but is
not the lead. The exposure-only benchmark reproduces the label at
least as well as the raw score on both ROC-AUC and PR-AUC. The
score's isolated contribution is the nested DeLong increment
($+\valExpIncrement$) and the within-opportunity-stratum AUC
($\valExpWithinAUC$). Combined exposure-plus-score logits (ROC-AUC
near $0.98$) are deliberately omitted from this table because
$E_i$ mechanically encodes the label and their fit is not score
evidence.
\end{tablenotes}
\end{threeparttable}
\end{adjustbox}
\end{table}

\begin{figure}[htbp]
\centering
\includegraphics[width=0.82\textwidth]{./output/figures/fig_observed_vs_expected_contact_bins}
\caption[Observed versus expected defendant contact by score bin]{Opportunity-adjusted
diagnostic: observed defendant contact $O_i$ versus expected
contact $E_i$ under random matching, by award-layer score bin.
The positive gap $O_i-E_i$ in the top score bins is the
within-opportunity signal; expected contact is built from
procurement opportunity cells (item, buyer, year, modality).
\textit{Alt text:} A binned plot with award-layer score bins on
the horizontal axis and defendant-contact counts on the vertical
axis. Two series are shown per bin, observed contact and expected
contact under random matching; observed contact exceeds expected
contact in the highest score bins, indicating excess contact
concentrated where the score is high.}
\label{fig:observed_vs_expected_contact}
\end{figure}

\subsection{Timing, Leakage, and Case Composition}
\label{sec:validation_timing}

The previous subsection isolated a small, exposure-adjusted
residual ordering. This subsection asks whether that ordering can
be deployed, and the honest answer is a downgrade. The ranking is
informative as a retrospective, incumbent-pool triage diagnostic,
but it is concentrated in one adjudicated case and one procurement
environment, and it does not support prospective, platform-wide
deployment without first handling new entrants, score-zero ties,
and cross-case breadth.

A retrospective rank measures; it does not deploy. Every benchmark
anchors its labels to CADE adjudications that close years after the
bidding they describe, so a high in-sample AUC certifies construct
validity, not field performance. The strict-timing design freezes
the score, the always-loser status, and the frequent-loser
threshold to the $\valYearStart$--$2016$ window and evaluates
against $2017$--$2019$ labels. There the signal survives \emph{only}
inside the always-loser training pool: the frequent-loser flag
attains ROC-AUC $\valStrictFLAUC$ and the continuous score
$\valStrictContAUC$. On the realistic full candidate universe
---every firm scorable at the screening date, with entrants carried
at score zero---discrimination collapses to ROC-AUC
$\valStrictFullAUC$, PR-AUC $0.005$, and \textbf{precision@500
$=\valStrictFullPrec$ and recall@500 $=\valStrictFullPrec$}. A
screen that returns zero true positives in its top $500$ candidates
is not deployable. The mechanical cause is structural: of the $193$
positives, $45$ ($\valEntrantPosShare$) are entrants with no
pre-window history, so the frozen score cannot rank them, only
$\valStrictRankablePos$ positives are rankable at all, and roughly a
quarter of the universe is tied at score zero. A loser-side score
built from past losing cannot order firms that have not yet lost.
The rolling-origin design confirms this: on the full universe every
test year from $2014$ to $2019$ returns PR-AUC near $0.004$--$0.006$
and precision@500 and recall@500 of zero, with the worst year
($2015$) at ROC-AUC $\valRollWorstAUC$, below chance; the apparent
within-pool improvement is retrospective accumulation of adjudicated
labels, not forward prediction. The rolling-origin curves are in
Appendix~\ref{app:validation_audits_submission}.

The decisive test is case composition. A single case
(trens/metros, CADE process $08700.004617/2013\text{-}41$) supplies
$\valTopCasePosShare$ of all positives and $\valTopCaseTP$ of the
top-$500$ true positives. Dropping it collapses the operational
metrics: PR-AUC falls from $\valLOCOfullPRAUC$ to
$\valLOCOdropLargestPRAUC$ ($-71\%$), precision@500 from
$\valFullPrec$ to $\valLOCOdropLargestPrec$, and recall@500 from
$0.342$ to $0.233$; dropping the top two cases leaves PR-AUC
$0.029$. ROC barely moves ($0.939\to0.929$), which is precisely why
ROC is the wrong lead metric here---it is insensitive to the case
concentration the operational metrics expose. Case-balanced
precision@500 is $\valCaseBalancedPrec$ against the firm-weighted
pooled $\valFullPrec$; only six linkable cases contribute any
top-$500$ true positives, the largest accounting for $64$--$87\%$.
The same concentration appears across environments: dropping the
largest item-group ($\valTopItemGroupShare$ of positives) roughly
halves PR-AUC ($0.126\to0.062$), and $96.8\%$ of positives sit in
Pregão. A clustered randomization-inference test
(buyer$\times$item-group$\times$year cells, $B=1000$) confirms the
ordering is real but thin: observed PR-AUC $\valLOCOfullPRAUC$ beats
the null at $p=\valClusterRIp$, yet case-coverage breadth is
\emph{not} significant ($p=\valClusterRIcovP$). The leave-one-case-out
distribution is reported in
Appendix~\ref{app:validation_audits_submission};
Table~\ref{tab:timing_case_holdout} synthesizes the full battery.

The most this design can support, even under strict timing, is a
triage queue. A retrospective, incumbent-pool, single-case ordering
tells an agency where to look first among firms it already observes;
it cannot identify cartels ex ante, rank firms it has never seen, or
be deployed platform-wide on the strength of one adjudicated
cluster. The composite verdict is C (the operational signal is
mainly retrospective and incumbent-pool) and D (it is case-dominant
and cluster-sensitive).

\begin{table}[htbp]
\centering
\caption{Timing and Case-Composition Synthesis}
\label{tab:timing_case_holdout}
\scriptsize
\begin{adjustbox}{max width=\textwidth,center}
\begin{threeparttable}
\begin{tabular}{p{2.7cm}p{2.6cm}ccccccp{1.9cm}p{2.6cm}}
\toprule
Design & Information set & $N$ & Pos. & ROC-AUC & PR-AUC & Prec@500 & Rec@500 & Leakage status & Interpretation \\
\midrule
Full-sample retrospective diagnostic & Full retrospective & $16{,}843$ & $\valExpNpos$ & $\valExpRawScoreAUC$ & $\valLOCOfullPRAUC$ & $\valFullPrec$ & $0.342$ & Upper bound; all leaks & Construct validity, not timing \\
Relaxed temporal holdout & Score frozen; future wins used & $16{,}843$ & $193$ & $\valAUCFLfirmTemp$ & --- & --- & --- & Leaks future wins & Score-timing sensitivity only \\
Strict full test universe & Strict; entrants at score 0 & $41{,}444$ & $193$ & $\valStrictFullAUC$ & $0.005$ & $\valStrictFullPrec$ & $\valStrictFullPrec$ & No leak & Collapse: not deployable \\
Strict rankable (incumbent) sample & Strict; prior-loss firms & $21{,}819$ & $\valExpNpos$ & $\valStrictFLAUC$ & $0.045$ & $0.062$ & $0.161$ & No leak & Survives only inside pool \\
Rolling-origin mean & Strict; full universe & --- & --- & $0.56$ & $0.005$ & $\valStrictFullPrec$ & $\valStrictFullPrec$ & No leak & Accumulation, not prediction \\
Rolling-origin worst year (2015) & Strict; full universe & $29{,}804$ & $127$ & $\valRollWorstAUC$ & $0.004$ & $\valStrictFullPrec$ & $\valStrictFullPrec$ & No leak & Below chance \\
Leave-one-case-out mean & Full retrospective & $16{,}750$ & --- & $0.95$ & --- & $\valCaseBalancedPrec$ & $0.522$ & Case-structured & Case-balanced precision tiny \\
Leave-largest-case-out & Full retrospective & $16{,}750$ & $86$ & $0.929$ & $\valLOCOdropLargestPRAUC$ & $\valLOCOdropLargestPrec$ & $0.233$ & Case-structured & $-71\%$ PR-AUC collapse \\
Leave-top-two-cases-out & Full retrospective & $16{,}649$ & $75$ & $0.925$ & $0.029$ & $0.030$ & $0.200$ & Case-structured & Thin residual tail \\
Clustered randomization inference & Label shuffle in cells & $16{,}843$ & $\valExpNpos$ & $0.939$ & $\valLOCOfullPRAUC$ & --- & --- & Permutation null & $p=\valClusterRIp$; coverage $p=\valClusterRIcovP$ \\
Direct-defendant scope & Strict and full scores & --- & --- & $\valDirectStrictAUC$ & --- & --- & --- & No leak & Scope boundary, not a detector \\
\bottomrule
\end{tabular}
\begin{tablenotes}
\scriptsize
\item \textit{Notes:} PR-AUC, precision@500, and recall@500 are
the primary metrics because the target is rare; ROC-AUC is
reported alongside but is not the lead, and it is deliberately
robust precisely because it is insensitive to the case
concentration that the operational metrics expose. The
full-sample retrospective row is a diagnostic \emph{upper bound},
not a timing claim. In the strict and rolling-origin designs the
score is frozen before the test window; later CADE adjudications
are still used as labels but were not legally observable at the
screening date, so these designs distinguish prospective scoring
from retrospective adjudication. Cobidders are
adjudication-anchored exposure labels, not cartel members. The
strict full universe carries entrants at score zero; a quarter of
that universe is tied at zero. ``$\valStrictFullPrec$'' denotes an
exact zero true positives in the top $500$.
\end{tablenotes}
\end{threeparttable}
\end{adjustbox}
\end{table}

\subsection{Direct-Defendant Scope and Bounded Reading}
\label{sec:validation_direct}
\label{sec:validation_summary}

The direct-defendant exercise aligns the validation target with the
screen's legal-economic role. Direct CADE defendants are the legal
anchors that locate adjudicated environments; they need not be
persistent zero-win firms, since defendants may occupy winner-side,
allocation, or rotating roles. The award-layer statistic ranks the
loser side by construction, so the empirical prediction is
asymmetry, and the data follow it: against direct defendants AUC is
$\valAUCdirectStd$, close to random. A generic cartel-membership
detector should rank defendants as well as cobidders; this one does
not, which is the expected result for a loser-side ranking. A
universe-anchored scope matrix sharpens the point---raw
participation count on the full BEC universe against direct
defendants returns AUC $\valAUCParticDefendants$, \emph{below}
random, so a high-participation firm is less likely to be a direct
defendant under the loser-side logic, reversing the direction
rather than merely failing. This is consistent with the role
separation in D4 (\valDirectShareAL\% of direct defendants are
always-losers; median win rate $\valDirectMedWR$ versus
$\valOthersMedWR$ for cobidders). The full eight-row scope matrix is
in Appendix~\ref{app:validation_audits_submission}.

Taken together, the benchmarks draw a boundary rather than certify a
screen. Each addresses a distinct pre-registered threat---legal
anchoring, participation volume, market exposure, timing, leakage,
and scope asymmetry---and the consolidated threat map is in
Appendix~\ref{app:validation_audits_submission}. Persistent
zero-win participation orders adjudication-anchored loser-side
exposure, and the ordering survives participation-volume and
opportunity-set discipline as a significant residual. But that
residual is modest, largely retrospective, concentrated among
incumbents and in one adjudicated case, and silent on direct
defendants. The defensible reading is therefore bounded: the award
layer ranks where loser-side exposure is most likely worth
bid-layer follow-up, within a reach the tests make explicit; it does
not validate the legal status of ranked firms, nor support
prospective, platform-wide deployment on its own. The next step is
economic interpretation: among firms already selected by the
frequent-loser rule, how far do cobidders display a distinct
loser-side profile rather than a mechanical collection of
high-volume losers concentrated in a few environments?
